\begin{definition}[The eight native checkerboard boundary registers]%
    \label{def:peps_kitaev_native_boundary_registers}
    \lean{TNLean.PEPS.kitaevNativeBoundaryEdge,
        TNLean.PEPS.kitaevNativeBoundaryEdgeEquiv,
        TNLean.PEPS.kitaevNativeBoundaryEquiv}
    \leanok
    \uses{def:peps_kitaev_native_checkerboard_block}
    For the locally oriented plaquette at $v=(x,y)$, enumerate its crossing
    bonds by a side $i\in\{N,E,S,W\}$ and a place $b\in\{0,1\}$ in that side.
    The places follow clockwise boundary order. This is a bijective
    enumeration of the actual crossing edges. It gives an equivalence
    $\mathcal E_v$ from native boundary configurations to four binary pairs.
    These are the boundary registers of the blocking construction in
    \cite[Section~7, ``Kitaev's code state'']{Schuch2010PEPS}.
\end{definition}

\begin{theorem}[The native boundary enumeration and its label reads]%
    \label{thm:peps_kitaev_native_boundary_registers}
    \lean{TNLean.PEPS.kitaevNativeBoundaryEdge_val,
        TNLean.PEPS.kitaevNativeBoundaryEquiv_apply}
    \leanok
    \uses{def:peps_kitaev_native_boundary_registers}
    Write $R_{a,b}$ for the actual right edge from $(a,b)$ and $U_{a,b}$
    for its up edge; these denote unordered native bonds even at a seam.
    The enumeration is
    \begin{align}
        N&=(U_{x,y+1},U_{x+1,y+1}),&
        E&=(R_{x+1,y+1},R_{x+1,y}),\\
        S&=(U_{x+1,y-1},U_{x,y-1}),&
        W&=(R_{x-1,y},R_{x-1,y+1}).
    \end{align}
    The equivalence $\mathcal E_v$ reads precisely the eight labels used
    in the native coefficient identity. Every native boundary assignment
    therefore occurs exactly once among these four pairs.
\end{theorem}

\begin{proof}\leanok
    Each bond is determined by its label read in the native block.
    The displayed formulas follow by evaluating these reads on indicator
    assignments. The right and up edge families are injective and disjoint;
    the period bounds also keep opposite boundary bonds distinct.
    Translate the plaquette to the coordinate rectangle of side lengths two.
    Its native perimeter contains eight crossing bonds. The eight distinct
    listed bonds consequently exhaust its boundary, and translation carries
    this enumeration to every position, including both seams.
\end{proof}

\begin{definition}[The native boundary CNOT and tensor maps]%
    \label{def:peps_kitaev_native_boundary_cnot}
    \lean{TNLean.PEPS.kitaevNativeBoundaryCNOT,
        TNLean.PEPS.kitaevNativeCheckerboardMatrix,
        TNLean.PEPS.kitaevNativeColorMatrix,
        TNLean.PEPS.kitaevNativeBoundaryZeroEmbedding}
    \leanok
    \uses{def:peps_kitaev_native_boundary_registers,
        def:peps_kitaev_checkerboard_maps}
    Let $C$ be the four pair CNOTs. Their native boundary action is
    $C_v=\mathcal E_v^{-1}C\mathcal E_v$.
    Let $\mathsf A_v$ be the actual open tensor map, with the four original
    native spins as its rows and all native crossing assignments as its columns.
    Let $\mathsf K_v$ be the color-difference tensor with the same physical
    rows, and let $E_v$ include the four colors into the native boundary
    registers by setting the four second registers to zero.
\end{definition}

\begin{theorem}[The exact native tensor-map identity after boundary CNOTs]%
    \label{thm:peps_kitaev_native_boundary_cnot}
    \lean{TNLean.PEPS.kitaevNativeBoundaryCNOT_pairs,
        TNLean.PEPS.kitaevNativeBoundaryCNOT_permMatrix_mem_unitaryGroup,
        TNLean.PEPS.kitaevNativeBoundaryZeroEmbedding_isIsometry,
        TNLean.PEPS.kitaevNativeCheckerboardMatrix_mul_boundaryCNOT}
    \leanok
    \uses{def:peps_kitaev_native_boundary_cnot,
        thm:peps_kitaev_native_boundary_registers,
        thm:peps_kitaev_native_checkerboard_identification,
        thm:peps_kitaev_checkerboard_map_identity}
    Reading the output of $C_v$ gives the four pair CNOTs. Its native
    permutation matrix $\mathsf C_v$ is unitary on the whole boundary space,
    and $E_v^\dagger E_v=1$. The actual native tensor map obeys
    \begin{align}
        \mathsf A_v\mathsf C_v=\mathsf K_v E_v^\dagger.
    \end{align}
    This is a local native tensor-map equality for periods at least three,
    including seam-crossing blocks. It does not assert a physical
    renormalization operation or a globally tiled checkerboard state.
\end{theorem}

\begin{proof}\leanok
    The derived boundary equivalence conjugates the pair CNOTs to a
    permutation of all native boundary configurations. Permutation matrices
    are unitary, and fixing four zero registers gives an isometric inclusion.
    The native coefficient identity identifies $\mathsf A_v$ with the
    explicit blocked map after reindexing its physical rows and boundary
    columns. Reindex the already proved identity
    $\mathsf B\mathsf C=\mathsf K E^\dagger$ by these same equivalences.
    Finite matrix multiplication is preserved by the bijective change of
    boundary indices, giving the asserted native equality.
\end{proof}
